\subsection*{同一命令为什么会占两个位置}

重新从一条独立历史开始。三副本的共同应用边界为 40，账户 1 余额 10000，购买、订单、流水、身份绑定均为空，E1 有效。本节动作是 \code{purchase.\allowbreak create.\allowbreak v1}。S1 取 \code{(1,\allowbreak purchase.\allowbreak create.\allowbreak v1,\allowbreak E1,\allowbreak K1)}，输入为一份 \code{paper-001}、报价 Q7、CNY 3000 分。N1 在任期 7 把它写入 41，复制并提交，然后应用：分配 P1/O1/D1，余额变为 7000，保存原结果。

N1 在返回以前失效。N2 接任，客户端重发原 S1 和原输入。假设 N2 已经补齐、应用到 41，但它的命令接收层只负责把请求排进日志，就可能在位置 43 再保存一份同样命令，位置 42 是新任期的空命令。Raft 保证 41 和 43 各自只有一个确定内容。它们却是两个不同位置，应用若每次都减 3000，就会得到 4000~\cite{raft,ongaro}。

前一节建立了同一日志位置的跨副本一致，本节增加同一业务身份的跨位置恢复。状态机应用 43 时，先按 S1 读取已经保存的绑定，核对规范化输入，再返回 P1/O1/D1 和当次 7000。位置 43 正常进入已应用前缀，业务流水数量保持为一。

\CElevenFigure{ch11-original09}{同一 S1 位于两个已提交日志位置，第二次应用恢复首次结果}{fig:ch11-duplicate-positions}

在所有副本上保存相同错误也可能完全一致。三台机器一起把余额变成 4000，能够通过各副本状态相等的检查，却违反同一 S 至多一条购买扣款的业务条件。正确性要同时检查复制前缀与应用规格。

\subsection*{E1 已经存在，为何原文又谈客户端会话}

第三章已给购买操作登记 E1，并用许可把账户、动作、E1、K1 与固定截止联系起来。E1 决定旧身份的准入边界，完整 S 决定查找哪项购买。Raft 原文在讨论重试时，使用客户端编号加命令序列号标记命令。博士论文还把每个客户端已完成的序号与响应组织为会话状态~\cite{raft,ongaro}。这里的客户端是同一组命令身份的协议持有者。一个持久客户端代理可以保存这些身份，并在连接或服务进程更换后继续使用。

直接保留完整 S，是本书状态机最自然的方案。命令携带 S，状态机保存第三章已经建立的映射 $M(S)$，同一 E1 内各个 K 独立进入映射。原输入、真实购买编号和原结果都从这份映射恢复。新的复制协议改变 M 与业务数据怎样持久存在，身份规则从原购买延续。

若客户端库要求 \code{client\_id/seq}，可以把账户、动作和 E1 编码为协议客户端，把 K 编码为序号。任意字符串 K 需要额外的双向唯一映射 U，选用这种扩展时，U 与业务绑定共同复制和快照。主路径直接使用完整 S，不经过这一适配。第~\ref{section:12-3}~节再比较任意键、可核验序号键及额外适配各自的保存成本。

E1 的固定截止和持久关闭规则继续有效。缺失或关闭的 E1 通过原购买入口返回过期。第二部分已经受理的 C1/T1 则由服务端原映射继续恢复，关闭购买准入以后，平台结果核对和未完成任务各有自己的持久责任。Raft 学位论文里的活跃客户端保活方案与本书固定期限会话具有不同的时间规则。本章把资格状态纳入复制，具体关闭、判期和回收步骤留在第十二章~\cite{ongaro}。

\subsection*{应用一次命令时共同改变什么}

先把当前状态摆出来。W 保存业务数据，包括余额、购买、订单、流水以及确定性编号计数器。M 保存从完整 S 到规范化版本 V、原输入 X、真实 P 和原结果 $R(S)$ 的绑定。$\Gamma$ 保存 E1 等登记会话及其准入状态。j 是这份状态对应的已应用日志边界。把它们合成一个状态：
\[
  \Sigma_j=(W_j,M_j,\Gamma_j,j).
\]

下标 j 表示这些内容都来自应用完同一日志前缀后的状态。例如一条同参恢复只把应用位置向前推进，原业务数据、绑定和结果保持。一次会话关闭命令改变 $\Gamma$，随后旧入口命令根据这个已复制决定返回过期。

\CElevenFigure{ch11-full-state}{同一应用边界 j 把业务数据、操作绑定、准入状态连成完整状态}{fig:ch11-full-state}

命令包含完整 S、输入解释版本、原始可核验输入和报价依据。入口先校验认证、操作许可和结构，状态机再按确定性准入规则与当前 $\Gamma$ 仲裁。涉及截止时刻的决定要让各副本取得同一依据，本章的数值历史限定 E1 已有效登记且全过程处在有效期间。第十二章会建立可复制的关闭决定。

已有 S 的输入解释沿原 V，先比较原 X，再决定恢复原结果或报告本次冲突。新 S 才进入首次报价与余额判断。第一次绑定方案允许一个尚未发布的候选被放弃，随后另一份同 S 输入先进入已提交序列。服务以首次正式应用形成的绑定为准，之后的异参请求保留原绑定，返回冲突。这与第四、六章区分参与者输入前提与服务责任的做法一致。

\CElevenFigure{ch11-binding-branches}{会话检查以后，已有 S 恢复原结果，新 S 才进入首次业务判断}{fig:ch11-binding-branches}

下面给出本地简化状态机完整的应用轮廓。整项转换串行发布，暂存状态在全部检查成功以前不对查询开放。业务拒绝保存的是首次确定结论，结构错误和输入冲突属于当前请求的处理结果：

\begin{BookCode}
输入：下一条已提交日志项 e，当前完整状态 Σ_j
要求 e 接在当前可应用前缀之后，先取得独占应用权
若 e 为协议空命令：发布相同业务状态和新应用边界，结束
核验 e 的命令格式，取得完整 S 与对应输入版本
核验已登记会话及许可的准入决定
若会话关闭或缺失：
    将本次响应设为过期，业务内容保持
否则若 M 已有 S：
    按原 V 解释来件，比较全部原输入
    同参：将本次响应设为原 P 和完整 R(S)，业务内容保持
    异参：将本次响应设为输入冲突，原 M/W 保持
否则：
    按受控规则核对报价及其账户、商品、数量、币种和金额
    若报价无效：
        将本次响应设为输入错误，不绑定已成交结果
    否则：
        在暂存状态分配真实 P
        若余额小于原金额：
            保存该 S/P 的稳定余额不足结果，不增加 O/D
        否则：
            减去原金额，分配真实 O/D
            同时保存 P/O/D 的归属、金额和原成交快照
            从真实编号和当次余额构造完整成功 R(S)
        同时保存 M(S)
        将本次响应设为该 S/P 的完整结果
核对余额、流水、身份、引用和结果的不变量
共同发布更新后的 W/M/Γ，应用边界推进到 e 的位置
将本次响应交回请求等待者
\end{BookCode}

冲突、过期和无效报价都已经是该已提交命令的应用结论，它们同样走共同出口推进边界。若命令内容损坏到无法按约定解释，应用引擎停止该条恢复路径并报告错误。正常格式的业务拒绝与这种存储或协议错误分开处理。

如果采用内存状态机，发布后的状态可以在进程崩溃后丢失，但已提交持久日志和完整快照必须足以重建它。若应用状态另存到本地事务数据库，则业务变化、M 和对应应用边界共同提交。后一种实现中，记录 j 却漏写 M，或保存扣款却没有原结果，都会在重放时产生不同恢复行为。

为使原结果可逐字段核对，本地 S1 沿用第三章的完整成功结果。许可中的固定截止与 K1 一起进入保存内容。这里的 7000 来自 S1 应用时的余额，作为原成交快照保存在结果中：

\begin{BookCode}
{
  "idempotency_key": "K1",
  "purchase_id": "P1",
  "status": "succeeded",
  "order_id": "O1",
  "debit_id": "D1",
  "product_id": "paper-001",
  "quote_id": "Q7",
  "currency": "CNY",
  "amount_minor": 3000,
  "balance_after_minor": 7000,
  "retry_until": "2026-10-18T10:00:00Z"
}
\end{BookCode}

实际编号可以由复制状态中的递增计数器确定，也可以先在命令里携带经规范核验的候选，再由状态机形成绑定。三个副本各自调用随机编号生成器，会使它们对同一命令返回不同 O。各副本分别读商品现价也有同样问题：一台看到 Q7，另一台看到 Q8，就可能扣出不同金额。将需要仲裁的外部输入带入共同命令，并按同一规则核验，才能得到相同转换~\cite{state-machine-approach}。

\subsection*{最近一个结果为什么覆盖不了并发}

Raft 扩展版的简述为每个客户端保存最近序号和响应。相关文献把并发情况另行展开，要求保存一组序号与响应，并结合客户端反馈决定可回收的响应范围~\cite{raft,ongaro}。简化结构的成立前提要放在算法旁边。

先看足以使用最近一个结果的受限协议。客户端在收到序号 n 的完整结果以前，只能重发 n，收到后才允许开始 n+1。服务保存 n 的结果时，正常客户端尚无机会提出 n+1。客户端以后提出 n+1，表示它已经获知 n。即使网络里迟到的 n 副本再出现，服务也可以将它识别为旧序号，按公布的旧请求规则处理。若接口要求任意多年以后都能回放 n 的原结果，就要保留比单个最近结果更完整的恢复依据。

本书允许同一 E1 有多项独立购买。另一条独立历史从余额 10000、空映射开始，K1、K2 都购买 3000。若用序号键把它们写为 1、2，两项可以同时在途，2 先应用时得到当次 7000，1 却可能尚未到达。服务仅凭最大已见序号 2 把 1 当成旧请求，就会漏掉一项合法购买。

若反过来先应用 1，再应用 2，服务只保存 2 的响应，则 1 的返回丢失后重试，原响应已被覆盖。返回 2 的结果会把 S1 指向另一购买，拒绝 1 又无法兑现原结果恢复。单个最大序号无法同时表示哪些较小请求已经执行、哪些尚未到达，以及每项操作自己的结果。

\CElevenFigure{ch11-original11}{序号 2 先到时，序号 1 可以尚未执行，两项未决购买分别保存结果}{fig:ch11-concurrent-sequences}

因此本章直接按完整 S 保留所有具备恢复资格的绑定，M 分别保存原输入和原结果。2 先到只说明 2 已有记录，1 到达后独立判断。并发请求的首次效果按照提交日志的应用次序串行形成，原编号和当次余额也沿这个次序产生。

这个序列号用于区分同一协议客户端的不同操作。业务真有前驱条件，例如第 2 步只能在第 1 步以后进行，就另行检查前驱并保存等待责任。第十章连续增量的 h 用于判断对象变化的前后关系，这里的逐笔绑定则分别回答每项 S 是否执行及返回什么。

\subsection*{从复制前缀接到业务不变量}

第六章的本地余额等式在这个封闭模型中继续成立。初始余额为 $b_0=10000$，已提交购买扣款流水集合为 D，每条 d 的正金额为 $a(d)$，则当前余额满足：
\[
  b=b_0-\sum_{d\in D}a(d),\qquad b\ge 0.
\]

状态机第一次应用 S1 时，余额减 3000，D 同时增加一条 3000，等式两边共同变化。再次应用同 S1，读原 M，余额与 D 都保持。新 S2 是另一项合法购买，余额再从 7000 减为 4000。余额不足的首次结论保存拒绝 P，没有支出流水，余额保持。

还要保留关系不变量：每个 S 至多一份已发布绑定，同一购买成功时恰好关联一张 O 和一条 D，D 的归属与金额和 P/O 一致。每个同参恢复返回原 $R(S)$，每条流水又能回到完整 S。初态为空时这些关系成立。首次成功、首次拒绝、同参恢复、异参冲突各自保持它们，协议空命令保持业务状态，副本追赶只按同一序列重建既有转换。

\CElevenFigure{ch11-invariants}{从共同初态逐项检查允许转换：余额守恒、单一绑定与原结果保持一起成立}{fig:ch11-invariants}

Raft 的前缀安全提供了一致的转换输入，状态机的归纳论证提供每次转换的业务正确性。RIFL 对完成记录的要求把两者接到恢复上：完成记录与业务变化原子保存，具有相应持久性，并且重试能够找到它~\cite{rifl}。在这里，M 随 W 一起复制、快照和恢复，接班者才能在相同 S 上继续这些关系。
\FloatBarrier
